\subsection{五曲線の全特殊値排除と中山固有値数の上界}
\label{subsec:cubic-four-simple-final-curves}
\begin{theorem}\label{thm:cubic-four-simple-sector-empty}
強い内在的H3の二次包絡の中山作用は、相異なる固有値を四つ持たない。
従って定理\ref{thm:cubic-three-eigenvalue-sector-empty}と合わせ、
全H3の中山作用の相異なる固有値数は高々二つである。
単一固有値の場合は固有値1であるが、単冪部分が自明とはここでは結論しない。
\end{theorem}
\begin{proof}
命題\ref{prop:cubic-four-simple-five-curves}の五表示を全て排除すればよい。
その番号0,3,7,10,13について、係数体を $F=\mathbb Q(\zeta_L)$ とし、
前節の式を次のように一変数化する（$z=\zeta_L$）。
\[
\begin{array}{c|r|c|c}
\text{番号}&L&(t_1,t_2)\text{ または }t_1&\text{非零条件}\\\hline
0&10&(2t^2/(t^2+z),t)&t(t^2+z)\ne0\\
3&10&(t,-2z^2t/(t^2-z^2))&t(t^2-z^2)\ne0\\
7&6&t&t\ne0\\
10&14&(2z^3t/(t^2-z^2),t)&t(t^2-z^2)\ne0\\
13&18&(t,2t/(t^2+z^4-z))&t(t^2+z^4-z)\ne0
\end{array}
\]
各分母が零になる場合は元rank式の非零定数項に矛盾するため、これは必要曲線の全点を覆う。

以下の局所化計算は、一般点のrankから特殊値のrankを推論しない。
元tensorのmiddle行列 $H$ から独立な六列 $R$ を選び、
$R$ の非零六次minor $\Delta$ を記録する。
その六列を二次関係とする自由語の全関係挿入で次数3,4の商を計算し、
全pivotの分子と分母を記録する。
最後に中山固定三次部分の正常条件
$Gv=\mu_B^{-1}(v)G$ の行列を構成し、同じく全pivotの分子・分母を記録する。
係数の分母と上表の定義域多項式も加える。
これら非零一変数多項式の積を $P$ とすると、$F[t,P^{-1}]$ 上で
全行操作が定義され、$\Delta$ が単元なので選んだ六関係は元middle像の全体と一致する。
商と左右積、正常行列はこの局所化上で正しく底変更される。
五表示で次数商は20,35、正常行列はそれぞれ
$26\times4$, $23\times4$, $22\times4$, $17\times2$, $14\times2$ の列full rankとなる。
従って $P\ne0$ の全点で $W_1(B)=0$ である。

$P=0$ の点も全て扱う。
$F[t]$ で全pivot多項式を因数分解し、積を掛け戻して元多項式との一致を検算する。
各相異なる既約因子 $f$ について既約性を確認し、
$E_f=F[t]/(f)$ 上で元の $H$ から次数商と正常行列を\emph{新しく}構成する。
上表の定義域外なら捨ててよいが、定義域内では全て商次元20,35かつ正常核零となる。
一つの根を数値代入したのではなく、既約因子の全根を同時に含む体での計算であり、
任意の代数閉標数0体への埋込みでもrank条件が保たれる。
\[
\begin{array}{c|r|r|r}
\text{番号}&\text{全相異なる因子}&\text{定義域内の因子}&\text{その次数和（}F\text{上）}\\\hline
0&11&8&20\\
3&5&2&3\\
7&13&12&23\\
10&8&5&11\\
13&9&6&9
\end{array}
\]
これで局所化の外に残る全点も排除した。
全四単純スペクトルの有限被覆、全非零support層、全Smith有限成分、零次元層、
前節の有限点と本節の五曲線を通じて、四単純H3領域は空である。
三固有値領域は既証定理により空。
単一固有値が1という主張は命題\ref{prop:h3-spectral-filter}による。
\end{proof}
公開再現器は\\
\texttt{verify\_cubic\_four\_simple\_final\_curves.py}、証明書は\\
\texttt{cubic\_four\_simple\_final\_curves\_certificate.json}。
全pivot多項式と全例外因子の係数を数体のべき基底で保存し、
各因子の元source rank・商次元・正常rankを保持する。
この定理は全CY被覆、単一固有値の全単冪持上げ、全intrinsic群oid・成分・相対降下の完成を意味しない。
